% --- Archon physics formalization source begin ---
% archon:physics
% archon:covers PhyXMiniProblems/problem_phyx_mini_0569.lean
% archon:source-report reports/phyx_mini/problem_phyx_mini_0569.source.json

\chapter{Physics problem phyx\_mini\_0569}
\label{ch:PhyXMiniProblems_problem_phyx_mini_0569}

\paragraph{Problem source.}
\#\# Physical scenario

An x-ray photon with energy of 100 keV experiences a ‘head-on’ collision with an electron at rest.

\#\# Question

Using conservation of momentum and energy, find the recoil velocity of the electron.

\#\# Answer choices

- A: A: 0.55c

- B: B: 0.95c

- C: C: 0.32c

- D: D: 10.2c

\#\# Figure caption (auxiliary; use the image as primary evidence)

The image consists of two diagrams side by side, showing similar setups with axes and particles.

**Left Diagram:**
- There are horizontal and vertical axes labeled "x" and "y," respectively.
- An orange squiggly line represents a photon, positioned along the x-axis and labeled "photon."
- A red circle labeled "electron" is positioned on the y-axis at the origin.
- The photon appears to be moving towards the red circle.

**Right Diagram:**
- Similar to the left, the axes are labeled with "x" and "y."
- An orange squiggly line, representing a photon, is positioned along the x-axis and labeled "photon."
- The red circle, labeled "electron," is positioned slightly along the positive x-axis, indicating motion in that direction, with an arrow indicating the movement to the right.

Overall, the diagrams depict the interactions between a photon and an electron, showing differing positions or movements before and after interaction.

\#\# Dataset metadata

Quantum Phenomena; Physical Model Grounding Reasoning, Multi-Formula Reasoning

\paragraph{Recorded answer/context.}
C: C: 0.32c

\paragraph{Figure/image path.}
/root/proposal\_for\_physic/science-mango/phyx\_mini\_run/phyx\_data/test\_image/569.png

\paragraph{Formalization target.}
create a compiling Lean file with sorry bodies at `PhyXMiniProblems/problem\_phyx\_mini\_0569.lean`. The Lean declarations must preserve the physical quantities, dimensions or dimensional roles, figure labels, governing-law hypotheses, and final relation expressed by this problem.
Use Mathlib/Physlib names found through LeanExplore where available. If a physics API is missing, introduce faithful local abstractions rather than scalar placeholder aliases.

\begin{theorem}[Physics formalization target]
\label{thm:physics:phyx_mini_0569:target}
\lean{PhyXMiniProblems.ProblemPhyXMini0569.electronRecoilVelocity_is_choice_C}
\uses{def:physics:phyx-mini-0569:phyxminiproblems-problemphyxmini0569-photonelectronbackscattersetup, def:physics:phyx-mini-0569:phyxminiproblems-problemphyxmini0569-matchesproblemdata, def:physics:phyx-mini-0569:phyxminiproblems-problemphyxmini0569-matchesprimarybackscatterfigure, def:physics:phyx-mini-0569:phyxminiproblems-problemphyxmini0569-hasphysicalbackscatterparameters, def:physics:phyx-mini-0569:phyxminiproblems-problemphyxmini0569-satisfiesrelativisticbackscatterlaws, def:physics:phyx-mini-0569:phyxminiproblems-problemphyxmini0569-electronrecoilspeedfractionofc, def:physics:phyx-mini-0569:phyxminiproblems-problemphyxmini0569-answerchoice, def:physics:phyx-mini-0569:phyxminiproblems-problemphyxmini0569-displayedspeedfractionofc, def:physics:phyx-mini-0569:phyxminiproblems-problemphyxmini0569-roundstonearesthundredth, def:physics:phyx-mini-0569:phyxminiproblems-problemphyxmini0569-isclosestanswerchoice}
The assigned autoformalize agent should translate this physics problem into Lean declarations in the covered file, with theorem and lemma proof bodies written as `by sorry`.
\end{theorem}
\begin{proof}
This is an autoformalization task, not a proof task. Produce statements that can later be proved by the physics prover without weakening the physical model.
\end{proof}

% --- Archon declaration topology begin ---
\section{Declaration topology}
\begin{definition}[Plane Momentum]
  \label{def:physics:phyx-mini-0569:phyxminiproblems-problemphyxmini0569-planemomentum}
  \lean{PhyXMiniProblems.ProblemPhyXMini0569.PlaneMomentum}
  A dimensionful two-dimensional momentum in the plane of the figure
\end{definition}
\begin{proof}
  This is the declared model definition of Plane Momentum.
\end{proof}
\begin{definition}[Diagram Axis]
  \label{def:physics:phyx-mini-0569:phyxminiproblems-problemphyxmini0569-diagramaxis}
  \lean{PhyXMiniProblems.ProblemPhyXMini0569.DiagramAxis}
  The two coordinate axes printed in both panels of the primary image
\end{definition}
\begin{proof}
  This is the declared model definition of Diagram Axis.
\end{proof}
\begin{definition}[to Fin]
  \label{def:physics:phyx-mini-0569:phyxminiproblems-problemphyxmini0569-diagramaxis-tofin}
  \lean{PhyXMiniProblems.ProblemPhyXMini0569.DiagramAxis.toFin}
  \uses{def:physics:phyx-mini-0569:phyxminiproblems-problemphyxmini0569-diagramaxis}
  The Fin 2 coordinate corresponding to a printed axis label
\end{definition}
\begin{proof}
  This is the declared model definition of to Fin.
\end{proof}
\begin{definition}[energy In Joules]
  \label{def:physics:phyx-mini-0569:phyxminiproblems-problemphyxmini0569-energyinjoules}
  \lean{PhyXMiniProblems.ProblemPhyXMini0569.energyInJoules}
  Joule readout of a physical energy
\end{definition}
\begin{proof}
  This is the declared model definition of energy In Joules.
\end{proof}
\begin{definition}[energy In Kilo Electron Volts]
  \label{def:physics:phyx-mini-0569:phyxminiproblems-problemphyxmini0569-energyinkiloelectronvolts}
  \lean{PhyXMiniProblems.ProblemPhyXMini0569.energyInKiloElectronVolts}
  \uses{def:physics:phyx-mini-0569:phyxminiproblems-problemphyxmini0569-energyinjoules}
  Kilo-electron-volt readout of a physical energy
\end{definition}
\begin{proof}
  This is the declared model definition of energy In Kilo Electron Volts.
\end{proof}
\begin{definition}[momentum Component In SI]
  \label{def:physics:phyx-mini-0569:phyxminiproblems-problemphyxmini0569-momentumcomponentinsi}
  \lean{PhyXMiniProblems.ProblemPhyXMini0569.momentumComponentInSI}
  \uses{def:physics:phyx-mini-0569:phyxminiproblems-problemphyxmini0569-planemomentum, def:physics:phyx-mini-0569:phyxminiproblems-problemphyxmini0569-diagramaxis}
  One momentum component in coherent SI units kg m / s
\end{definition}
\begin{proof}
  This is the declared model definition of momentum Component In SI.
\end{proof}
\begin{definition}[momentum Magnitude In SI]
  \label{def:physics:phyx-mini-0569:phyxminiproblems-problemphyxmini0569-momentummagnitudeinsi}
  \lean{PhyXMiniProblems.ProblemPhyXMini0569.momentumMagnitudeInSI}
  \uses{def:physics:phyx-mini-0569:phyxminiproblems-problemphyxmini0569-planemomentum, def:physics:phyx-mini-0569:phyxminiproblems-problemphyxmini0569-momentumcomponentinsi}
  Euclidean magnitude of a plane momentum in kg m / s
\end{definition}
\begin{proof}
  This is the declared model definition of momentum Magnitude In SI.
\end{proof}
\begin{definition}[speed In Meters Per Second]
  \label{def:physics:phyx-mini-0569:phyxminiproblems-problemphyxmini0569-speedinmeterspersecond}
  \lean{PhyXMiniProblems.ProblemPhyXMini0569.speedInMetersPerSecond}
  Metres-per-second readout of a nonnegative physical speed
\end{definition}
\begin{proof}
  This is the declared model definition of speed In Meters Per Second.
\end{proof}
\begin{definition}[speed Of Light In Meters Per Second]
  \label{def:physics:phyx-mini-0569:phyxminiproblems-problemphyxmini0569-speedoflightinmeterspersecond}
  \lean{PhyXMiniProblems.ProblemPhyXMini0569.speedOfLightInMetersPerSecond}
  Physlib's exact vacuum speed of light, read in metres per second
\end{definition}
\begin{proof}
  This is the declared model definition of speed Of Light In Meters Per Second.
\end{proof}
\begin{definition}[Photon State]
  \label{def:physics:phyx-mini-0569:phyxminiproblems-problemphyxmini0569-photonstate}
  \lean{PhyXMiniProblems.ProblemPhyXMini0569.PhotonState}
  \uses{def:physics:phyx-mini-0569:phyxminiproblems-problemphyxmini0569-planemomentum}
  A photon state retains both its dimensionful energy and spatial momentum
\end{definition}
\begin{proof}
  This is the declared model definition of Photon State.
\end{proof}
\begin{definition}[Figure Panel]
  \label{def:physics:phyx-mini-0569:phyxminiproblems-problemphyxmini0569-figurepanel}
  \lean{PhyXMiniProblems.ProblemPhyXMini0569.FigurePanel}
  The two panels represented side by side in the primary image
\end{definition}
\begin{proof}
  This is the declared model definition of Figure Panel.
\end{proof}
\begin{definition}[Particle Label]
  \label{def:physics:phyx-mini-0569:phyxminiproblems-problemphyxmini0569-particlelabel}
  \lean{PhyXMiniProblems.ProblemPhyXMini0569.ParticleLabel}
  Particle labels printed in each panel
\end{definition}
\begin{proof}
  This is the declared model definition of Particle Label.
\end{proof}
\begin{definition}[Horizontal Motion]
  \label{def:physics:phyx-mini-0569:phyxminiproblems-problemphyxmini0569-horizontalmotion}
  \lean{PhyXMiniProblems.ProblemPhyXMini0569.HorizontalMotion}
  Qualitative arrow directions visible in the supplied raster
\end{definition}
\begin{proof}
  This is the declared model definition of Horizontal Motion.
\end{proof}
\begin{definition}[Photon Electron Backscatter Figure]
  \label{def:physics:phyx-mini-0569:phyxminiproblems-problemphyxmini0569-photonelectronbackscatterfigure}
  \lean{PhyXMiniProblems.ProblemPhyXMini0569.PhotonElectronBackscatterFigure}
  \uses{def:physics:phyx-mini-0569:phyxminiproblems-problemphyxmini0569-diagramaxis, def:physics:phyx-mini-0569:phyxminiproblems-problemphyxmini0569-figurepanel, def:physics:phyx-mini-0569:phyxminiproblems-problemphyxmini0569-particlelabel, def:physics:phyx-mini-0569:phyxminiproblems-problemphyxmini0569-horizontalmotion}
  Presentation-level information transcribed from the two-panel image. It contains no numerical speed, energy, or answer-choice value
\end{definition}
\begin{proof}
  This is the declared model definition of Photon Electron Backscatter Figure.
\end{proof}
\begin{definition}[Photon Electron Backscatter Setup]
  \label{def:physics:phyx-mini-0569:phyxminiproblems-problemphyxmini0569-photonelectronbackscattersetup}
  \lean{PhyXMiniProblems.ProblemPhyXMini0569.PhotonElectronBackscatterSetup}
  \uses{def:physics:phyx-mini-0569:phyxminiproblems-problemphyxmini0569-planemomentum, def:physics:phyx-mini-0569:phyxminiproblems-problemphyxmini0569-photonstate, def:physics:phyx-mini-0569:phyxminiproblems-problemphyxmini0569-photonelectronbackscatterfigure}
  Independent physical quantities before and after the collision. The recoil speed is an unconstrained dimensionful field here; it is related to the electron's energy and momentum only by the governing law below
\end{definition}
\begin{proof}
  This is the declared model definition of Photon Electron Backscatter Setup.
\end{proof}
\begin{definition}[Matches Problem Data]
  \label{def:physics:phyx-mini-0569:phyxminiproblems-problemphyxmini0569-matchesproblemdata}
  \lean{PhyXMiniProblems.ProblemPhyXMini0569.MatchesProblemData}
  \uses{def:physics:phyx-mini-0569:phyxminiproblems-problemphyxmini0569-energyinkiloelectronvolts, def:physics:phyx-mini-0569:phyxminiproblems-problemphyxmini0569-photonelectronbackscattersetup}
  Numerical data stated or standard for this instance. The incident photon has energy 100 keV; the standard electron rest energy is calibrated as 511 keV. Neither field specifies the outgoing speed
\end{definition}
\begin{proof}
  This is the declared model definition of Matches Problem Data.
\end{proof}
\begin{definition}[Matches Primary Backscatter Figure]
  \label{def:physics:phyx-mini-0569:phyxminiproblems-problemphyxmini0569-matchesprimarybackscatterfigure}
  \lean{PhyXMiniProblems.ProblemPhyXMini0569.MatchesPrimaryBackscatterFigure}
  \uses{def:physics:phyx-mini-0569:phyxminiproblems-problemphyxmini0569-momentumcomponentinsi, def:physics:phyx-mini-0569:phyxminiproblems-problemphyxmini0569-momentummagnitudeinsi, def:physics:phyx-mini-0569:phyxminiproblems-problemphyxmini0569-photonelectronbackscattersetup}
  Qualitative and directional evidence from the primary image. In particular, the outgoing photon reverses direction while the electron recoils to the right; every physical momentum has zero transverse component
\end{definition}
\begin{proof}
  This is the declared model definition of Matches Primary Backscatter Figure.
\end{proof}
\begin{definition}[Has Physical Backscatter Parameters]
  \label{def:physics:phyx-mini-0569:phyxminiproblems-problemphyxmini0569-hasphysicalbackscatterparameters}
  \lean{PhyXMiniProblems.ProblemPhyXMini0569.HasPhysicalBackscatterParameters}
  \uses{def:physics:phyx-mini-0569:phyxminiproblems-problemphyxmini0569-energyinjoules, def:physics:phyx-mini-0569:phyxminiproblems-problemphyxmini0569-speedinmeterspersecond, def:physics:phyx-mini-0569:phyxminiproblems-problemphyxmini0569-speedoflightinmeterspersecond, def:physics:phyx-mini-0569:phyxminiproblems-problemphyxmini0569-photonelectronbackscattersetup}
  Positivity and subluminality conditions for the physical event
\end{definition}
\begin{proof}
  This is the declared model definition of Has Physical Backscatter Parameters.
\end{proof}
\begin{definition}[Satisfies Relativistic Backscatter Laws]
  \label{def:physics:phyx-mini-0569:phyxminiproblems-problemphyxmini0569-satisfiesrelativisticbackscatterlaws}
  \lean{PhyXMiniProblems.ProblemPhyXMini0569.SatisfiesRelativisticBackscatterLaws}
  \uses{def:physics:phyx-mini-0569:phyxminiproblems-problemphyxmini0569-diagramaxis, def:physics:phyx-mini-0569:phyxminiproblems-problemphyxmini0569-energyinjoules, def:physics:phyx-mini-0569:phyxminiproblems-problemphyxmini0569-momentumcomponentinsi, def:physics:phyx-mini-0569:phyxminiproblems-problemphyxmini0569-momentummagnitudeinsi, def:physics:phyx-mini-0569:phyxminiproblems-problemphyxmini0569-speedinmeterspersecond, def:physics:phyx-mini-0569:phyxminiproblems-problemphyxmini0569-speedoflightinmeterspersecond, def:physics:phyx-mini-0569:phyxminiproblems-problemphyxmini0569-photonelectronbackscattersetup}
  Governing relations for an isolated relativistic photon--electron collision. The two photon laws are E = pc; the outgoing electron obeys E² = (m c²)² + (pc)² and v = p c² / E. These general relations do not contain the requested numerical recoil speed
\end{definition}
\begin{proof}
  This is the declared model definition of Satisfies Relativistic Backscatter Laws.
\end{proof}
\begin{definition}[electron Recoil Speed Fraction Of C]
  \label{def:physics:phyx-mini-0569:phyxminiproblems-problemphyxmini0569-electronrecoilspeedfractionofc}
  \lean{PhyXMiniProblems.ProblemPhyXMini0569.electronRecoilSpeedFractionOfC}
  \uses{def:physics:phyx-mini-0569:phyxminiproblems-problemphyxmini0569-speedinmeterspersecond, def:physics:phyx-mini-0569:phyxminiproblems-problemphyxmini0569-speedoflightinmeterspersecond, def:physics:phyx-mini-0569:phyxminiproblems-problemphyxmini0569-photonelectronbackscattersetup}
  The requested dimensionless electron recoil speed v/c
\end{definition}
\begin{proof}
  This is the declared model definition of electron Recoil Speed Fraction Of C.
\end{proof}
\begin{lemma}[electron Recoil Speed Fraction from conservation]
  \label{lem:physics:phyx-mini-0569:phyxminiproblems-problemphyxmini0569-electronrecoilspeedfraction-from-conservation}
  \lean{PhyXMiniProblems.ProblemPhyXMini0569.electronRecoilSpeedFraction_from_conservation}
  \uses{def:physics:phyx-mini-0569:phyxminiproblems-problemphyxmini0569-energyinjoules, def:physics:phyx-mini-0569:phyxminiproblems-problemphyxmini0569-photonelectronbackscattersetup, def:physics:phyx-mini-0569:phyxminiproblems-problemphyxmini0569-matchesprimarybackscatterfigure, def:physics:phyx-mini-0569:phyxminiproblems-problemphyxmini0569-hasphysicalbackscatterparameters, def:physics:phyx-mini-0569:phyxminiproblems-problemphyxmini0569-satisfiesrelativisticbackscatterlaws, def:physics:phyx-mini-0569:phyxminiproblems-problemphyxmini0569-electronrecoilspeedfractionofc}
  Conservation and the relativistic dispersion laws determine the exact recoil ratio for an arbitrary incident photon energy E and electron rest energy M = m\_e c² in the head-on backscatter geometry: v/c = 2 E (E + M) / (M² + 2 M E + 2 E²)
\end{lemma}
\begin{proof}
  The conclusion follows from the governing relations and figure readouts named in the statement of electron Recoil Speed Fraction from conservation.
\end{proof}
\begin{definition}[Answer Choice]
  \label{def:physics:phyx-mini-0569:phyxminiproblems-problemphyxmini0569-answerchoice}
  \lean{PhyXMiniProblems.ProblemPhyXMini0569.AnswerChoice}
  Labels of the four speed choices displayed by the source
\end{definition}
\begin{proof}
  This is the declared model definition of Answer Choice.
\end{proof}
\begin{definition}[displayed Speed Fraction Of C]
  \label{def:physics:phyx-mini-0569:phyxminiproblems-problemphyxmini0569-displayedspeedfractionofc}
  \lean{PhyXMiniProblems.ProblemPhyXMini0569.displayedSpeedFractionOfC}
  \uses{def:physics:phyx-mini-0569:phyxminiproblems-problemphyxmini0569-answerchoice}
  The numerical coefficient of c printed beside each answer label
\end{definition}
\begin{proof}
  This is the declared model definition of displayed Speed Fraction Of C.
\end{proof}
\begin{definition}[recorded Dataset Answer]
  \label{def:physics:phyx-mini-0569:phyxminiproblems-problemphyxmini0569-recordeddatasetanswer}
  \lean{PhyXMiniProblems.ProblemPhyXMini0569.recordedDatasetAnswer}
  \uses{def:physics:phyx-mini-0569:phyxminiproblems-problemphyxmini0569-answerchoice}
  Answer label recorded in the source dataset
\end{definition}
\begin{proof}
  This is the declared model definition of recorded Dataset Answer.
\end{proof}
\begin{definition}[Rounds To Nearest Hundredth]
  \label{def:physics:phyx-mini-0569:phyxminiproblems-problemphyxmini0569-roundstonearesthundredth}
  \lean{PhyXMiniProblems.ProblemPhyXMini0569.RoundsToNearestHundredth}
  Agreement with a speed coefficient displayed to the nearest hundredth
\end{definition}
\begin{proof}
  This is the declared model definition of Rounds To Nearest Hundredth.
\end{proof}
\begin{definition}[Is Closest Answer Choice]
  \label{def:physics:phyx-mini-0569:phyxminiproblems-problemphyxmini0569-isclosestanswerchoice}
  \lean{PhyXMiniProblems.ProblemPhyXMini0569.IsClosestAnswerChoice}
  \uses{def:physics:phyx-mini-0569:phyxminiproblems-problemphyxmini0569-answerchoice, def:physics:phyx-mini-0569:phyxminiproblems-problemphyxmini0569-displayedspeedfractionofc}
  A displayed choice is at least as close to the modeled ratio as all others
\end{definition}
\begin{proof}
  This is the declared model definition of Is Closest Answer Choice.
\end{proof}
% --- Archon declaration topology end ---
% --- Archon physics formalization source end ---
